% ======================================================================
% app-datalog-proofs.tex -- Datalog^neg: soundness/completeness of the
% layered choice basis, the full negation-semantics proof, and the
% support-level monus elimination proof.
% ======================================================================
\section{Datalog with Negation: Proofs}
\label{sec:appendix-datalog}

This appendix proves the claims of Section~\ref{sec:datalog}: soundness and
completeness of the canonical layered choice basis, the negation-semantics
correspondence, and support-level monus elimination.

\subsection{Soundness and Completeness of the Choice Basis}

\begin{proposition}[Soundness and completeness]
  \label{prop:datalog-sc}
  For finite $\mbox{Datalog}^\neg$ programs whose negative SCCs admit the
  layered decomposition (each SCC has finitely many local stable extensions
  given its fixed context, and composing them in topological order is sound
  and complete for global stable models), the canonical layered choice basis
  of Definition~\ref{def:datalog-canonical} is sound and complete: each
  resolving determination yields a stable model, and every stable model is
  yielded by exactly one resolving determination.
\end{proposition}
\begin{proof}
  The Gelfond--Lifschitz reduct~\cite{gelfond1988stable} decomposes along the
  SCC DAG: once atoms outside $C_i$ are fixed, the reduct of $C_i$'s rules
  depends only on those atoms. \emph{Soundness:} composing local stable
  extensions in topological order yields a global stable model, since each
  local extension is a minimal model of its local reduct and their composition
  is a minimal model of the full reduct. \emph{Completeness:} given a stable
  model $M$, define $D_M$ by the shared sealing prefix followed by the choice
  $\varphi_{C_i} = M|_{C_i}$ at each layer $k{+}i$; local minimality of
  $M|_{C_i}$ follows from global minimality of $M$. The map $M \mapsto D_M$ is
  injective, giving uniqueness.
\end{proof}

\subsection{Negation Semantics as Filtration Levels (Full Proof)}

\begin{proof}[Proof of Theorem~\ref{thm:negation-general}]
  \emph{Depth $k{+}d$.} The $k$ sealing layers are shared by all resolving
  determinations (stratified evaluation is unique), so they add no cell
  distinction. Each choice layer $k{+}i$ resolves $C_i$ using atoms fixed by
  layers $k{+}1, \dots, k{+}i{-}1$; since $C_i$'s rules reference earlier
  cycles, those layers do not commute, giving depth $k{+}d$.

  \emph{(a) Stable models.} Each resolving determination discharges all $k{+}d$
  layers, producing a two-valued assignment that is a stable model
  (Proposition~\ref{prop:datalog-sc}); conversely every stable model is one
  such determination. Hence stable models biject with $\Dets$, and
  stable-model semantics reads $\mathcal{F}_{k+d} = \mathcal{P}(\Dets)$.

  \emph{(b) Well-founded semantics.} WFS applies the alternating
  fixpoint~\cite{vanGelder1991wellFounded} after the sealing prefix. At choice
  layer $k{+}i$ the alternating fixpoint checks whether $C_i$ has a unique
  stable extension given the discharged prefix; if so the layer is
  \emph{forced} and WFS discharges it, otherwise $C_i$'s atoms are classified
  $\mathbf{u}$ and no dependent layer is discharged. Let $d^\ast$ be the number
  of consecutive forced layers from $k{+}1$. WFS reads
  $\mathcal{F}_{k+d^\ast}$: atoms settled by layers $1, \dots, k{+}d^\ast$ have
  full or empty support ($\mathbf{t}/\mathbf{f}$), and atoms depending on
  layers beyond $k{+}d^\ast$ have partial support ($\mathbf{u}$).

  \emph{(c) Stratified semantics.} When $d = 0$ the sealing prefix resolves
  all atoms, $|\Dets| = 1$, and $\mathcal{F}_k = \{\emptyset, \Dets\}$.
\end{proof}

\noindent
\emph{Existence of stable models.} The correspondence presumes stable models
exist; programs such as $a \leftarrow \neg a$ have none, and are excluded by
the decomposition hypothesis (an SCC with no local stable extension yields
$\Dets = \emptyset$, where every $\mathcal{F}_k = \{\emptyset\}$). The
WFS$\leftrightarrow$stable-support step uses the Van Gelder--Ross--Schlipf
characterization of the well-founded model as the intersection structure of
stable models~\cite{vanGelder1991wellFounded}.

\subsection{Monus Elimination (Full Proof)}

Let $P$ be finite and stratified with strata $S_0, \dots, S_k$, and let $K$ be
naturally ordered, zero-divisor-free, and $\omega$-continuous. Both the
sealing and monus~\cite{dannert2021semiring} strategies evaluate bottom-up by
least fixpoint at each stratum over $(K,+,\cdot)$, agreeing on positive rules;
they differ only on a negated literal $\neg b$ (with $b$ already evaluated to
$\mathrm{val}(b) \in K$): sealing contributes $1_K$ if $\mathrm{val}(b) = 0_K$
and $0_K$ otherwise, while monus contributes
$1_K \mathbin{\dot{-}} \mathrm{val}(b)$.

\begin{lemma}
  \label{lem:monus}
  In such $K$, $1_K \mathbin{\dot{-}} v = 1_K$ if $v = 0_K$ and
  $1_K \mathbin{\dot{-}} v = 0_K$ if $v \neq 0_K$.
\end{lemma}
\begin{proof}
  If $v = 0_K$ then $\max\{c \mid c \le 1_K\} = 1_K$. If $v \neq 0_K$, suppose
  $c \neq 0_K$ with $v + c \le 1_K$. Monotonicity gives $vc \le v, c$, so
  $vc + vc \le v + c \le 1_K$, and zero-divisor-freeness gives $vc \neq 0_K$.
  Repeated squaring yields nonzero $(vc)^{2^n}$ all satisfying
  $w + w \le 1_K$; by $\omega$-continuity their infimum $e$ is nonzero,
  idempotent, with $e + e \le 1_K$. Writing $e + f = 1_K$: if $f \neq 0_K$
  then $ef \neq 0_K$ and $ef + ef \le 1_K$, a nonzero element below $e$---%
  contradicting infimality---so $f = 0_K$ and $e = 1_K$. But then
  $1_K + 1_K \le 1_K$ forces $1_K \le 0_K$, a contradiction.
\end{proof}

\begin{proof}[Proof of Theorem~\ref{thm:monus-elimination}]
  By Lemma~\ref{lem:monus} the two strategies assign the same zero/nonzero
  value to every negated literal. Positive evaluation is identical, and $K$ is
  zero-divisor-free (a product is nonzero iff all factors are) and zerosumfree
  (a sum is nonzero iff some summand is), so induction on strata gives that
  every atom has the same zero/nonzero annotation under sealing and monus.
  Supports therefore agree.
\end{proof}

\subsection{$K$-Provenance for Negation-Licensed Atoms}

\begin{proof}[Proof of Proposition~\ref{prop:kprov-negation}]
  Let $M = \rho(D)$. By definition of a stable
  model~\cite{gelfond1988stable}, the Gelfond--Lifschitz reduct $P^{M}$---obtained
  by deleting every rule whose body contains a negative literal $\neg\beta$ with
  $\beta \in M$, and striking the remaining (satisfied) negative literals from
  the surviving rules---is a \emph{negation-free} Datalog program, and $M$ is
  its least model. Thus $P^{M}$ is positive and $M = P^{M}{\uparrow}\omega$, the
  least fixpoint of its immediate-consequence operator.

  Positive Datalog admits $K$-relational provenance in the sense of Green et
  al.~\cite{green2007provenance}: annotate each base fact by an indeterminate
  (or an application-supplied element of $K$), evaluate the least fixpoint with
  $+$ for alternative derivations and $\cdot$ for joint ones, and read the
  annotation $\Prov_K(\operatorname{Real}(M), a_{Q,\alpha})$ of a query output
  $\alpha$. Set $\operatorname{Real}(M)$ to this $K$-annotated least model of
  $P^{M}$. Soundness and completeness of $K$-provenance for positive programs
  give, under positivity of $K$, that the annotation is nonzero iff $\alpha$ is
  derivable in $P^{M}$, i.e.\ iff $\alpha \in M$ (and, for a query $Q$, iff
  $\alpha \in Q(M)$). This is precisely adequacy
  (Definition~\ref{def:consumer}) for the result-extensional consumer with
  $\Cond = \operatorname{Real}\circ\rho$, so the consumer of
  Example~\ref{ex:gkt} instantiated at $P^{M}$ is adequate.

  Crucially, an atom $\alpha$ whose derivation in $P^{M}$ uses a rule that
  originally carried a negative literal $\neg\beta$ (with $\beta \notin M$) is
  handled no differently: in $P^{M}$ that literal has been struck, so the rule
  is an ordinary positive rule and contributes to $\alpha$'s annotation by the
  same $+/\cdot$ rules. Hence negation-licensed atoms carry the same
  $K$-provenance as any other derived atom. Finally, since $\mathbb{N}[X]$ is
  the free commutative semiring on the base annotations, every commutative
  semiring $K'$ is the image of a unique semiring homomorphism
  $h:\mathbb{N}[X]\to K'$, and $h$ commutes with least-fixpoint
  evaluation~\cite{green2007provenance}; applying $h$ to
  $P_D(a_{Q,\alpha})\in\mathbb{N}[X]$ yields the corresponding specialized
  annotation (counting in $\mathbb{N}$, probability via a tuple-independent PDB
  homomorphism, access control, tropical cost, and so on), each sound for the
  negation-licensed atoms as well.
\end{proof}

\noindent
Proposition~\ref{prop:kprov-negation} is a level-$0$ (within-determination)
statement; it says nothing about which determinations exist or how an atom's
support varies across them. The cross-model structure is supplied separately by
the support filtration (Section~\ref{sec:filtration}), and the two are composed
by Theorem~\ref{thm:filtration-respected}.
